\documentclass[11pt,a4paper]{article}
\usepackage{fontspec}
\setmainfont{DejaVu Sans}[Scale=0.90]
\usepackage{amsmath,amssymb}
\usepackage[margin=2.2cm]{geometry}
\usepackage{booktabs,array,tabularx}
\usepackage{enumitem}
\usepackage{titlesec}
\usepackage{parskip}
\titlelabel{\thetitle.\quad}
\titleformat{\section}{\normalfont\large\bfseries}{\thesection.}{0.6em}{}
\titleformat{\subsection}{\normalfont\normalsize\bfseries}{\thesubsection}{0.6em}{}
\titlespacing*{\section}{0pt}{15pt}{5pt}
\titlespacing*{\subsection}{0pt}{11pt}{4pt}
\setlist[itemize]{leftmargin=1.4em,itemsep=1pt,topsep=3pt}
% Preserve fonts/margins; permit a final line-breaking pass for long prose.
\setlength{\emergencystretch}{2em}
\begin{document}
\begin{center}
{\Large\bfseries Fermeture cinétique : dérivation complète}\\[4pt]
{\normalsize Modèle gris isotrope, conservation et convention GK historique}\\[2pt]
{\small Partie I --- 16 septembre 2026}
\end{center}

\section{Résultat et domaine de validité}
On achève ici la fermeture \emph{linéaire, grise et isotrope} annoncée dans la note~09.
On ne dérive pas un modèle spectral de l'AlN réel. Le calcul ci-dessous est explicite :
distribution déplacée, moments angulaires, projection déviatorique, élimination du
moment rapide, puis réduction au modèle de face avant.

Une distinction est nécessaire. Avec la température définie par l'énergie locale,
la fermeture conservatrice donne
\[
 \tau_R\partial_t\mathbf q+\mathbf q=-\lambda\nabla T+
 \ell^2\left[\nabla^2\mathbf q+\frac13\nabla(\nabla\cdot\mathbf q)\right],
 \qquad
 \lambda=\frac13Cv^2\tau_R,\quad
 \ell^2=\frac15v^2\tau_N\tau_R .
\]
Les coefficients microscopiques sont ici donnés à l'ordre dominant en
$\tau_N/\tau_R$. Le coefficient $1/3$ diffère du $2$ de l'équation GK
historique employée jusqu'ici. Cette différence n'est pas une découverte du
projet : voir Sendra et al. (2022), équation~(21), et Struchtrup (2025).
La dérivation indépendante qui suit indique exactement où intervient la trace.

Hypothèses : cristal uniforme, vitesse $v$ commune aux modes, dispersion
$\omega_{\mathbf k}=v|\mathbf k|$, temps $\tau_N,\tau_R$ constants et positifs,
petites perturbations autour de $T_0$, collisions conservant l'énergie,
pas de source volumique. Les temps résistifs sont ceux du \emph{volume}
(umklapp, défauts, isotopes dans l'approximation grise). Les frontières demandent
des conditions de bord ; on ne les compte pas une seconde fois dans $\tau_R$.
L'élimination du moment rapide exige $|p|\tau_N\ll1$ et
$v\tau_N/L\ll1$ ; la limite collective ajoute $\tau_N/\tau_R\ll1$.
$L$ est l'échelle des variations, qui peut être plus petite que l'épaisseur.

\section{Équilibre déplacé et compatibilité}
Posons $\epsilon_{\mathbf k}=\hbar v|\mathbf k|$ et
$d\Gamma=d^3k/(2\pi)^3$, avec multiplicité de branches identiques incluse
dans les intégrales. La distribution immobile et sa version déplacée sont
\[
 f^0(T)=\frac1{\exp(\epsilon_{\mathbf k}/k_BT)-1},\qquad
 f^d(T,\mathbf u_d)=
 \frac1{\exp[(\epsilon_{\mathbf k}-\hbar\mathbf k\cdot\mathbf u_d)/(k_BT)]-1}.
\]
À l'ordre linéaire, pour $\theta=T-T_0$ et $\mathbf n=\mathbf k/|\mathbf k|$,
\[
 f^d-f^0(T_0)=
 \partial_T f^0(T_0)\left(\theta+\frac{T_0}{v}\mathbf n\cdot\mathbf u_d\right).
\]
En effet $\partial_T f^0=f^0(1+f^0)\epsilon/(k_BT_0^2)$.
On définit $E=\int\epsilon f\,d\Gamma$ et
$C=\partial E^0/\partial T|_{T_0}$. La température énergétique impose
$E-E_0=C\theta$. Au premier ordre, le terme de dérive est impair et ne change
pas $E$. La contrainte sur le quasi-moment fixe $\mathbf u_d$.

Le modèle de Callaway est
\[
 \partial_t f+v n_i\partial_i f
 =-\frac{f-f^d}{\tau_N}-\frac{f-f^0}{\tau_R}.
\]
Les états cibles doivent satisfaire
$\int\epsilon(f-f^d)d\Gamma=\int\epsilon(f-f^0)d\Gamma=0$ et
$\int\hbar k_i(f-f^d)d\Gamma=0$.
Ces égalités sont des \emph{conditions de compatibilité}, pas une conséquence
automatique de la présence de deux termes de relaxation.
Pour des temps dépendant du mode, les contraintes portent sur les intégrales
pondérées par $1/\tau_{\mathbf k}$ ; les étapes grises suivantes ne se transposent
pas en remplaçant simplement chaque temps par une moyenne arithmétique.

\section{Intégrales angulaires et premiers moments}
Avec $\langle\cdot\rangle=(4\pi)^{-1}\int_{S^2}\cdot\,d\Omega$,
l'isotropie et $n_i n_i=1$ donnent
\[
 \langle n_i n_j\rangle=\frac{\delta_{ij}}3,\qquad
 \langle n_i n_j n_k n_l\rangle=
 \frac{\delta_{ij}\delta_{kl}+\delta_{ik}\delta_{jl}
       +\delta_{il}\delta_{jk}}{15}.
\]
Pour établir $1/15$, on écrit la forme isotrope avec un coefficient $A$,
puis contracte $i=j$ et $k=l$ : $1=A(9+3+3)$.
Les moments impairs sont nuls par inversion de $\mathbf n$.

Introduisons une densité énergétique angulaire $g$, normalisée par
$\langle g\rangle=E-E_0$ et $\mathbf q=v\langle\mathbf n g\rangle$.
Sa partie déplacée linéarisée est
\[
 g^d=C\theta+\frac3v\,\mathbf n\cdot\mathbf q,\qquad
 \mathbf q=\frac13CT_0\mathbf u_d .
\]
Les moments d'ordre deux et trois (perturbations) sont
\[
 Q_{ij}=v^2\langle n_i n_j g\rangle,\qquad
 R_{ijk}=v^3\langle n_i n_j n_k g\rangle.
\]
La partie déplacée donne
\[
 Q^d_{ij}=\frac{v^2C\theta}{3}\delta_{ij},\qquad
 R^d_{ijk}=\frac{v^2}{5}
 (q_i\delta_{jk}+q_j\delta_{ik}+q_k\delta_{ij}).
\]
Pour \emph{toute} distribution, $Q_{ii}=v^2C\theta$.
On doit donc écrire $Q_{ij}=v^2C\theta\,\delta_{ij}/3+\Pi_{ij}$
avec $\Pi_{ii}=0$.
Enfin $\mathbf q=v^2\mathbf P$ pour
$\mathbf P=\int\hbar\mathbf k f\,d\Gamma$.
Cette identité exige la même vitesse linéaire isotrope : elle n'est pas
une loi universelle pour toutes les branches d'un cristal.

\section{Bilans exacts dans le modèle gris}
Les moments d'ordre zéro et un de l'équation cinétique donnent
\[
 C\partial_t\theta+\partial_iq_i=0,\qquad
 \partial_tq_i+\frac{Cv^2}{3}\partial_i\theta+\partial_j\Pi_{ij}
 =-\frac{q_i}{\tau_R}.
\]
Le terme normal disparaît du bilan de flux grâce à la compatibilité du
quasi-moment. Le moment d'ordre deux donne
\[
 \partial_t Q_{ij}+\partial_k R_{ijk}
 =-\frac{\Pi_{ij}}{\tau_c},\qquad
 \tau_c=\left(\tau_N^{-1}+\tau_R^{-1}\right)^{-1}.
\]
La trace reproduit le bilan d'énergie, puisque $R_{iik}=v^2q_k$.
La projection déviatorique est donc obligatoire.

À l'ordre nécessaire à l'élimination du moment rapide, on remplace
$R$ par $R^d$. Avec le bilan d'énergie,
\[
 \partial_t\Pi_{ij}+
 \frac{v^2}{5}\left(\partial_iq_j+\partial_jq_i
             -\frac23\delta_{ij}\partial_kq_k\right)
 =-\frac{\Pi_{ij}}{\tau_c}.
\]
Le terme $-\frac23\delta_{ij}\nabla\cdot\mathbf q$ résulte de
$-\frac{Cv^2}{3}\delta_{ij}\partial_t\theta$ après projection ;
il garantit une trace nulle.

\section{Élimination du moment rapide}
Si $\tau_c\partial_t\Pi$ est petit devant $\Pi$,
\[
 \Pi_{ij}=-\frac{v^2\tau_c}{5}
 \left(\partial_iq_j+\partial_jq_i
             -\frac23\delta_{ij}\partial_kq_k\right).
\]
Cette relation est une fermeture au premier ordre, pas une identité exacte
à toute fréquence. La contribution du moment d'ordre trois non déplacé et
la mémoire de $\Pi$ appartiennent aux termes négligés.
Prendre la divergence donne
\[
 \partial_j\Pi_{ij}
 =-\frac{v^2\tau_c}{5}
   \left(\nabla^2q_i+\frac13\partial_i\nabla\cdot\mathbf q\right).
\]
Substituer dans le bilan de flux et multiplier par $\tau_R$ produit
\[
 \tau_R\partial_t\mathbf q+\mathbf q
 =-\lambda\nabla T+\ell_c^2
   \left[\nabla^2\mathbf q+\frac13\nabla(\nabla\cdot\mathbf q)\right],
 \quad \lambda=\frac13Cv^2\tau_R,\quad
 \ell_c^2=\frac15v^2\tau_R\tau_c .
\]
Garder $\tau_c$ décrit la fermeture quasi stationnaire du moment d'ordre deux.
Cela ne prouve pas l'exactitude de tous les autres termes à ordre supérieur.
Dans la limite collective, $\tau_c=\tau_N[1+O(\tau_N/\tau_R)]$ et
$\ell_c^2\simeq\ell^2=v^2\tau_R\tau_N/5$.

\subsection{Origine de la forme historique}
Si l'on supprime $\partial_tQ_{ij}$ \emph{avant} de séparer sa trace énergétique,
on obtiendrait à tort, pour la même température,
\[
 Q_{ij}\simeq\frac{Cv^2\theta}{3}\delta_{ij}
 -\frac{v^2\tau_c}{5}
   (\partial_iq_j+\partial_jq_i+\delta_{ij}\nabla\cdot\mathbf q).
\]
Sa trace vaut $Cv^2\theta-v^2\tau_c\nabla\cdot\mathbf q$ au lieu de
$Cv^2\theta$. Sa divergence produit précisément le coefficient historique $2$.
Cette opération peut appartenir à une autre approximation en temps lent,
mais ne constitue pas la fermeture conservatrice générale annoncée.
L'accord en régime stationnaire solénoïdal ne résout pas le désaccord transitoire.
La forme historique reste un modèle constitutif utilisable et étudié :
son interprétation cinétique doit être indiquée.

\section{Réduction 1D et raccord au code}
Écrivons de façon générale le terme spatial
$\ell^2[\nabla^2\mathbf q+\alpha\nabla(\nabla\cdot\mathbf q)]$.
Pour un champ strictement longitudinal $\mathbf q=q(z,t)\mathbf e_z$,
\[
 \tau_R q_t+q=-\lambda T_z+L^2 q_{zz},
 \qquad L^2=(1+\alpha)\ell^2,\qquad \tau_\ell=L^2/a,\quad a=\lambda/C.
\]
\begin{center}
\begin{tabular}{@{}llll@{}}
\toprule
Convention & $\alpha$ & $L^2$ & $\tau_\ell$ (limite collective)\\
\midrule
GK historique & $2$ & $3\ell^2$ & $9\tau_N/5$\\
Conservatrice grise & $1/3$ & $4\ell^2/3$ & $4\tau_N/5$\\
\bottomrule
\end{tabular}
\end{center}
Le code reçoit directement $\tau_\ell$ ; aucune formule du quadripôle ne change.
Il faut fournir $4\tau_N/5$ pour la seconde convention, et $9\tau_N/5$
seulement pour la première. Les longueurs physiques déduites d'un même
$\tau_\ell$ diffèrent : $\ell_{\rm CE}=(3/2)\ell_{\rm hist}$.

À conditions initiales nulles, le bilan d'énergie donne
$\phi'=-Cp\theta$, donc $\phi''=-Cp\theta'$.
La loi de flux devient
\[
 \phi=-\lambda_{\rm eff}(p)\theta',\qquad
 \lambda_{\rm eff}(p)=\lambda\frac{1+p\tau_\ell}{1+p\tau_R}.
\]
On obtient $\theta''=\sigma^2\theta$ avec
\[
 \sigma^2=\frac{p}{a}\frac{1+p\tau_R}{1+p\tau_\ell},\qquad
 M=\begin{pmatrix}
 \cosh(\sigma d)&\sinh(\sigma d)/(\lambda_{\rm eff}\sigma)\\
 \lambda_{\rm eff}\sigma\sinh(\sigma d)&\cosh(\sigma d)
 \end{pmatrix}.
\]
$M$ relie l'état arrière à l'état avant, flux positif vers $z>0$.
Les limites $p=0$ se prennent analytiquement.

\section{Résonance, entropie et portée expérimentale}
La résonance $\tau_\ell=\tau_R$ annule le facteur constitutif et redonne
Fourier pour ce problème aux limites sans source ni perturbation initiale.
Son existence algébrique ne dépend pas de $\alpha$.
La traduction $\tau_R/\tau_N=9/5$ appartient à la convention historique.
La conversion conservatrice dominante donnerait formellement $4/5$ ;
aucun de ces rapports n'est dans $\tau_N/\tau_R\ll1$.
Avec $\tau_c$ conservé, $\tau_\ell/\tau_R=
4\tau_N/[5(\tau_R+\tau_N)]<4/5$ : pas d'égalité à temps positifs
dans cette fermeture. Ce constat n'est pas une trajectoire en température.

Pour $L^2\ge0$, la production 1D au second ordre autour de $T_0$ est
\[
 s=s_{\rm eq}(E)-\frac{\tau_Rq^2}{2\lambda T_0^2},\quad
 J_s=\frac qT+\frac{L^2qq_z}{\lambda T_0^2},\quad
 \sigma_s=\frac{q^2+L^2q_z^2}{\lambda T_0^2}\ge0.
\]
L'entropie n'impose donc pas à elle seule la valeur cinétique de $\alpha$.
De même, une mesure 1D dépend de $L^2$ seulement : si $\ell$ est libre,
elle ne distingue pas $\alpha=2$ de $\alpha=1/3$ par simple ajustement.
Une géométrie multidimensionnelle et des informations constitutives indépendantes
sont nécessaires pour tester cette distinction.

\section{Contrôles et références}
Les tests angulaires du projet intègrent numériquement sur la sphère, vérifient
les moments d'ordres deux et quatre, la conservation de la correction,
la trace de $\Pi$ et le coefficient longitudinal. Un test de réponse vérifie
l'identité des deux paramétrisations à $L^2$ fixé.
Ils contrôlent cette dérivation grise ; ils ne valident pas le modèle de l'AlN.

\begin{enumerate}
\item J. Callaway, Phys. Rev. \textbf{113}, 1046 (1959) ;
R. A. Guyer et J. A. Krumhansl, Phys. Rev. \textbf{148}, 766 et 778 (1966) :
origine des modèles, et non reproduction complète de leurs opérateurs ici.
\item L. Sendra et al., Phys. Rev. B \textbf{106}, 155301 (2022),
doi:\allowbreak 10.1103/\allowbreak PhysRevB.106.155301.
Manuscrit intégral consulté ; équation (21), discussion du coefficient $1/3$.
\item H. Struchtrup, \emph{Phonon hydrodynamics revisited} (2025),
doi:\allowbreak 10.1007/\allowbreak 978-3-031-93918-1\_27.
Notice institutionnelle et résumé consultés ; le chapitre intégral n'a pas
été obtenu par le dépôt institutionnel.
\item M. G. Hennessy et T. G. Myers (2021),
doi:\allowbreak 10.1007/\allowbreak 978-3-030-64272-3\_2.
Manuscrit accepté intégral consulté ; forme historique, équation (2).
\end{enumerate}
\end{document}
